\chapter{The discovery: shared generation rules}
\label{ch:discovery}

\begin{saybox}[title=What to say at the start of this section]
This is the pivot of the whole meeting. Ask the question out loud before
giving the answer: ``when we score a held-out row, are we testing whether
the model can predict a \emph{new relationship}, or only whether it can
produce another sample of a relationship it has already learned?'' Then
show the measurement.
\end{saybox}

\section{The question that had not been asked}

The diagnostic phase eliminated the model, the partition sizes, the
validation design, the cross-validation count, and the specialist
granularity. What remained was the variable along which the data was cut.

The original protocol groups by \texttt{context\_id}. The relevant question
is whether the constraint that this enforces is the constraint that this
dataset requires. In a corpus of measurements, grouping by experimental
unit is exactly right. In a corpus of \emph{generated} rows, the unit that
carries the information is the generator, not the sample.

\section{Measuring the overlap}

The production splitter was re-executed with the production seed on each
of the six specialists, and the resulting partitions were audited at two
levels: context membership and rule membership
(\evd{scripts/experiments/rule\_overlap\_audit.py}).

\tbl{tab_rule_overlap}{Rule-level audit of the original protocol. Contexts
never cross the split; generation rules always do.}

\begin{findingbox}[title=The central measurement]
In all six specialists:
\begin{itemize}[leftmargin=1.2em, itemsep=1pt]
\item \textbf{Context overlap between train and test: exactly $0$.} The
  protocol enforces its intended constraint perfectly.
\item \textbf{Test rows originating from a generation rule that also
  appears in training: $100.0\%$.} Not a majority --- every single row.
  $1\,320/1\,320$ for soft/general, $720/720$ for soft/safety,
  $1\,080/1\,080$, $600/600$, $900/900$, $480/480$.
\end{itemize}

Every test row in every production specialist was generated by a rule the
model was trained on.
\end{findingbox}

This is not a subtle statistical artefact requiring a delicate argument.
It is a structural property of the corpus interacting with the grouping
variable, and it is measurable in a single pass over the data.

\section{Why it happens mechanically}

Each generation rule spans many contexts. Soft/general contains seven
rules spread across its contexts; the largest rule alone contributes
$5\,654$ rows. A random $70/15/15$ allocation of \emph{contexts} will,
with overwhelming probability, place at least one context from every
sizeable rule in each partition. For a rule spanning $m$ contexts, the
probability that none of them lands in the test partition is roughly
$0.85^{m}$, which is negligible for all but the smallest rules.

The protocol therefore cannot avoid rule overlap by chance; it would have
to exclude it by construction, and it does not attempt to.

\section{What the model is actually being asked to do}

Return to \Cref{eq:generator}: rule $g$ produces rows via
$y = \phi_g(\mathbf{x};\theta_g) + \varepsilon$. Under the original
protocol the training set contains many draws from rule $g$ and the test
set contains further draws from the same $g$. The model's task at test time
is to evaluate a function it has already estimated at a new input drawn
from the same input distribution.

This decomposes the apparent test score into two components that the
protocol cannot separate:
\begin{equation}
\underbrace{\rsq_{\text{context-holdout}}}_{\text{what was measured}}
\;=\;
h\Big(
\underbrace{\text{interpolation within known rules}}_{\text{dominant}},
\;
\underbrace{\text{transfer to unknown rules}}_{\text{untested}}
\Big).
\end{equation}
Because every test row comes from a known rule, the second component
receives zero weight. The measured quantity is the first component alone.

\begin{findingbox}[title=Why this also explains the negative results]
The generation-rule account is not an excuse for the high scores; it has
to explain the low ones too, and it does. Where a specialist's rules
differ sharply from one another in their target scale --- as
\Cref{ch:folds} shows for hard/safety-endpoint --- even interpolation
within known rules is hard, because the model must serve several
incompatible relationships simultaneously. High scores and negative scores
in the same table are both consistent with a corpus built from a small
number of heterogeneous mechanisms.
\end{findingbox}

\section{What this does \emph{not} establish}

Three careful non-claims, each of which the panel may probe:

\begin{enumerate}[leftmargin=1.5em, itemsep=2pt]
\item \textbf{It is not leakage in the technical sense.} No feature
  contains target information unavailable at prediction time.
  \texttt{source\_rule\_id} is excluded from the model inputs. The models
  never see the rule label; the \emph{split} sees it, or rather fails to.
\item \textbf{It is not classical overfitting.} The models do not memorise
  individual rows and fail on held-out ones; they generalise well within
  rules, which is a real, if narrow, capability.
\item \textbf{It does not make the original numbers arithmetically wrong.}
  They correctly report performance under the protocol that produced them.
  What changes is the interpretation that may be attached to them.
\end{enumerate}

\section{The consequence for evaluation design}

If the quantity of interest is the ability to predict a formulation whose
governing relationship was not in the training data --- which is the
quantity relevant to any deployment claim --- then the evaluation must
withhold generation rules, not contexts. That is the subject of the next
chapter.

\fig{fig19_loro_diagram}{0.95}{%
  The revised protocol. Each fold withholds one entire generation rule
  from training and scores the model on that rule's rows alone.}
